\section{Exact disk encoding and the loop-counting law}\label{sec:discrete}

The purpose of this section is to obtain the exact disk laws used in the
later peeling comparison.  We first express the sphere law through
fully packed triangle disks.  Their partition functions then give both
the volume local limit theorem and the counting intensity of a marked
sphere loop.  All perimeter counts refer to boundary occurrences.

\subsection{The finite ensemble and its Tutte encoding}\label{disk:encoding}

We retain throughout the exact weight
\[
 4^{\ell(M,A)/2}=2^{\ell(M,A)}=2^{2k(A)+|A|-|V(M)|}.
\]
In particular the vertex-count factor remains in every decomposition below.
The triangles introduced here are combinatorial encoding cells; the conformal surface and graph distance remain those specified in Theorem~\ref{thm:main}.

Put \(X=(4\sqrt{2})^{-1}\), \(c_*=\sqrt 8\), \(B_n=(2/\pi)\sqrt n\log n\).
In the radial quadrangulation (corners joined by links from the face center to the primal vertex), split each quadrangle by the occupied edge if present and by the complementary dual edge otherwise, both drawn through the edge midpoint.
Call the two resulting triangles per edge Tutte triangles.
Color the primal vertices blue and face centers red.
The triangle has a monochromatic edge (its base) and two incidence links.
The flag interface in Theorem~\ref{thm:main} goes exactly once through each Tutte triangle from link to link; in the case of blue base it cuts across at face-center coordinate \(1/3\) in the flags, in the case of red base at primal-vertex coordinate \(2/3\).
This is the usual fully packed interface model.
The resulting disjoint simple loops have total count \(\ell(M,A)\) (complementary regions correspond to primal and dual components; disjoint circles on a sphere), and color the vertices according to nesting parity.

\begin{proposition}[Partition function and loop-counting intensity]\label{disk:palm}

For a loop let \(k,l\) be the numbers of its visited Tutte triangles with blue, respectively red base.
Call the side adjacent to the blue bases (i.e. toward the blue endpoints of crossed links) its blue side, and let \(u\) be the \(m_n\)-area of that side of the loop on the sphere, including everything nested there.
Write \(d_s(v)=s^2 v^{-2}\exp(-s^2/v)\) for \(s,v>0\), extended by zero at \(v=0\).
Then
\begin{equation}
\begin{gathered}
X^{2n}Z_n(4)\sim \frac{8\pi^2}{n^2\log^2 n},\\
\lim_n \mathbb E\sum_{\rm loops}\psi(k/B_n,u)
 =\frac1{2\pi^2}\int_0^\infty\int_0^1
 \psi(s,v)\frac{ds\,dv}{s^5}\,d_s(v)d_s(1-v).
\end{gathered}
 \label{eq:original-3}
\end{equation}
Here \(\psi\) is continuous with compact support in \((0,\infty)\times[0,1]\), extended by zero at length zero.
This does not assert that loop lengths can be read from a geometric limit or that every macroscopic loop has \(k\) in a range of order \(B_n\).

\end{proposition}

The scale $B_n$ in this proposition is a perimeter-count scale.  Its proof occupies the remainder of this section.

\subsection{Disk partition functions and exact peeling}\label{disk:peeling}

Use the fully packed disk model with root face of degree \(p\), all other faces triangles, with disjoint simple loops in the dual crossing all triangles from edge to edge (degree two per internal dual vertex, degree zero at root face, including incidence multiplicities).
Weight \(2^{\#\mathrm{loops}}x^{N}\), where \(N\) is the number of internal triangles.
The root directed edge keeps the root face on its right; the boundary need not be simple, and we use just the singleton point for \(p=0\).
Write \(F_p(x)\) for the total.
One fixes the boundary color, giving a unique coloring by the loops, with either choice having this same partition function.
Conversely paired triangles across monochromatic bases in the boundaryless, colored sphere model are exactly quadrangles: remove the bases, keeping their colors as diagonal choices. (Each face has exactly one base incidence, a crossed edge changes vertex color, and all gluing is by occurrences.) This recovers a bipartite quadrangulation and by drawing its blue diagonals the unique decorated primal map.
These descriptions via corners/edge sides allow bridges and loops on all maps.

\begin{lemma}[Slitting and the gasket]\label{disk:slit-gasket}
The degree-two disk partition function satisfies Equation~\eqref{eq:original-4}, and the gasket satisfies Equation~\eqref{eq:original-5} with the face weights given below.  Given the exposed boundary lengths and edge occurrences, every unexplored disk has its fresh Boltzmann law.
\end{lemma}
\begin{proof}

Slitting an oriented occupied primal root edge to make a root face of degree two, with the same orientation on the copy with slit on its right, does not alter the number of triangles or interfaces.
Even a self-edge gives a degree-two face by inserting two copies in the map's rotations; the boundary can pinch at the vertex.
Conversely a disk of degree two with internal triangles always has two distinct boundary edges (if both boundary incidences are the same bridge no boundary length is left for either attachment, hence no triangles).
With boundary blue, close the degree-two face to an occupied edge to undo the slit.
Since primal-dual exchange preserving weight gives exactly half the weight to occupied roots, and the bridge-only disk has weight one,
\begin{equation}
F_2(x)=1+\frac12\sum_{n\ge1}Z_n(4)x^{2n}.
 \label{eq:original-4}
\end{equation}

The outer gasket, obtained by retaining edges accessible without crossing a loop, has the ordinary boundary-map face weights
\[
 g_j(x)=2\sum_{l\ge0}\binom{j+l-1}{l}x^{j+l}F_l(x)\qquad(j\ge1).
\]
Indeed each internal face of the gasket is replaced with a triangle ring carrying an outermost loop and then an independent disk of opposite boundary color.
With a specified first outer base, the ring is determined by its interleaved cyclic list of \(j\) outer and \(l\) inner bases.
Radial links between successive triangles encode the boundary correspondence, with a rooted inner disk read starting at a prescribed inner base (or a single point if no such base); corners of faces and triangle ring boundaries need not have distinct labels in the assembled map.
Deletion/retention of all the edges not belonging to the ring gives the inverse, including in pinched cases.
In particular peeling the root edge of the gasket satisfies
\begin{equation}
F_p=\sum_{i=0}^{p-2} F_i F_{p-2-i}+\sum_{j\ge1}g_j F_{p+j-2}.
 \label{eq:original-5}
\end{equation}
The first step either deletes a root bridge incident to the root face on both sides or deletes a root edge merging an internal face with the root face.
Each exposed disk (and each loop interior) has fresh Boltzmann law given its length, by the factorization of the weights with distinguished boundary edge occurrences.

\end{proof}

\subsection{Resolvent and the critical singularity}\label{disk:resolvent}

We use two generating-function inputs: the critical disk asymptotic of
Da Silva--Hu--Powell--Wong, Theorem~1.4, and the one-cut resolvent of
Borot--Bouttier--Guitter \citep[Sections~6.1--6.3]{BBGBending}, recalled in
\citep[Equations~(2.20)--(2.23) and~(3.4)--(3.5)]{DSHPW}.
In our units the first input says
\[
 F_p(X)<\infty,\qquad
 f_p:=c_*^{-p}F_p(X)\sim\frac8{\pi^2}\frac{\log p}{p^2}.
\]
For $0<x<X$, the second input gives
\[
 F_p(x)=\int_a^b y^p\rho(y)\,dy,
 \qquad a\leq0<b,\quad |a|\leq b,
\]
where $\rho$ is a probability density, positive and continuous in the cut
interior and vanishing at its endpoints.  In the interior it satisfies
\[
 {\rm PV}\int_a^b\left(\frac1{y-w}+\frac1{2c-y-w}\right)\rho(w)\,dw=\frac{y}{2},\qquad c=\frac1{2x}.
\]
Here the perimeter model is exactly the fully packed triangulation model with symmetric face weight \(x\), loop fugacity \(2\), and no other weights as above.
Admissibility in the one-cut theorem requires the pointed partition function to be finite.  At fixed perimeter \(p\), a vertex mark on the gasket is also a vertex mark on the underlying triangulated disk, so it costs at most \(C_p(N+1)\) when there are \(N\) internal triangles.  For \(x<X\), the bound
\[
 \sup_{N\ge0}(N+1)(x/X)^N<\infty
\]
and \(F_p(X)<\infty\) therefore give the required pointed finiteness.
Also \(b\le c_*\) by the moment formula and monotonicity (use even moments), hence \(b<c\); the face series in the one-cut equation is regular at the cut.
This public input is about disk generating functions, not the requested joint geometric convergence \citep[Theorem~1.4 and Sections~2.4--3.1]{DSHPW}.

Put \(\alpha=c-b,\ \beta=c-a,\ A=\alpha^2,\ B=\beta^2\), and write \(r(t)=\rho(c-\sqrt t)\).
Then
\[
 {\rm PV}\int_A^B\frac{r(s)}{s-t}\,ds=\frac{c-\sqrt t}{2},\qquad
 \int_A^B\frac{r(t)}{2\sqrt t}\,dt=1.
\]

\begin{lemma}[Hilbert inversion]\label{disk:hilbert}

Let \(U^2\) have the arcsine probability law on \([A,B]\) with \(U>0\), and \(m_j=\mathbb E[U^j]\), \(d(t)=\sqrt{(t-A)(B-t)}\).
Hilbert inversion gives
\begin{equation}
m_1=c,\qquad r(t)=\frac{d(t)}{2\pi}\,\mathbb E\frac1{U+\sqrt t}.
 \label{eq:original-6}
\end{equation}

\end{lemma}\begin{proof}

Indeed divide \(\Phi(z)=\int r(s)/(z-s)\,ds\) by \(D(z)=\sqrt{(z-A)(z-B)}\sim z\); the jump from lower to upper value (upper minus lower) of that quotient is \((\sqrt t-c)/(i d(t))\), so
\[
 \Phi(z)/D(z)=\tfrac12\,\mathbb E[(U-c)/(z-U^2)].
\]
Here one can apply Cauchy's formula deforming onto the cut (Sokhotski--Plemelj in the regular interior; equivalently distributional boundary values).
Endpoint circle integrals vanish by boundedness of \(r\): the integral of \(|\Phi|\) along an endpoint circle of radius \(\epsilon\) is \(O(\epsilon\log(1/\epsilon))\) by integrating the absolute kernel first.
Decay at infinity gives \(m_1=c\), and taking the jump of \(\Phi\), using the vanishing interior Hilbert transform of the arcsine law, yields \eqref{eq:original-6}.
For independent \(U,V\) with this same law put \(h(U)=-U^3+(A+B)U-AB/U\).
Normalization in \eqref{eq:original-6} now reads \(1=\mathbb E[h(U)/(U+V)]/4\); symmetrizing and using \(m_2=(A+B)/2\) gives
\[
 c^2-AB m_{-1}^2=8.
\]
It follows as \(x\uparrow X\) that \(\alpha\to0\), \(\beta\to\pi c_*/2\) (by \(m_1=c\)).
In particular the left endpoint is strictly away from \(-c_*\), including for the limiting cut of the moment formula at \(X\), obtained by continuity or dominated convergence in \eqref{eq:original-6}.

\end{proof}

For precision set \(z=x^2/X^2,\ \sigma=A/B\), and write
\[
 J_\pm(\sigma)=\frac2\pi\int_0^{\pi/2}(\sigma\cos^2\theta+\sin^2\theta)^{\pm1/2}d\theta.
\]
The equations read
\begin{equation}
\beta=c/J_+,\qquad 1-z=\sigma(J_-/J_+)^2.
 \label{eq:original-7}
\end{equation}

\begin{lemma}[Singularity and coefficient asymptotics]\label{disk:singularity}

Locally in slit sectors as \(\delta=1-z\to0\),
\[
 \sigma\sim \frac{4\delta}{\log^2(1/\delta)},\qquad
 \frac{d\sigma}{d\delta}\sim\frac4{\log^2(1/\delta)}.
\]

Moreover, Equation~\eqref{eq:original-8} holds in these sectors, and the partition-function equivalent in Proposition~\ref{disk:palm} follows.
\end{lemma}\begin{proof}

We mean here a holomorphic continuation to small \(\delta\ne0,\ |\arg\delta|<\pi-\epsilon_0\) for each fixed \(\epsilon_0>0\), with uniform equivalents.
Indeed \(J_+\to2/\pi,\ J_+'=O(\log(1/|\sigma|)),\ J_-\sim\log(1/\sigma)/\pi,\ J_-'\sim-1/(\pi\sigma)\) uniformly in such sectors.
For example after \(v=\tan\theta\) the integral for \(J_-\) uses kernel \((\sigma+v^2)^{-1/2}(1+v^2)^{-1/2}\); after differentiation one uses \(\int_0^\infty(\sigma+v^2)^{-3/2}dv=1/\sigma\) and the sector bound by powers of \((|\sigma|+v^2)\).
Rouché on small fixed relative-radius circles about \(4\delta/\log^2(1/\delta)\) gives a unique solution there and the implicit derivative gives the second estimate.
This branch agrees by \eqref{eq:original-7} with the real solution near criticality.

For real \(c>c_*\) near \(c_*\) differentiate at fixed \(t\) in the cut:
\[
 \partial_c r(t)=\frac{t-\eta}{2\pi d(t)},\qquad
 \eta=\frac{m_1}{m_{-1}}=\frac{c^2}{J_+ J_-}.
\]
Indeed endpoints depend differentiably by \eqref{eq:original-7}, and \eqref{eq:original-6} lets us differentiate under the integrals (integrable endpoint inverse square-root behavior, smooth on compact interior).
Thus \(\partial_c\Phi\) has sum of its two boundary values \(-1\) there.
Multiplication of \(\partial_c\Phi+1/2\) by \(D\) removes this cut, leaving a polynomial by the endpoint bounds and the expansion at infinity, necessarily of the form \((z'-\eta)/2\) in the Stieltjes-transform coordinate \(z'\); normalization differentiated fixes \(\eta\).
In particular with primes in the next line meaning \(c\) derivatives of the moments \(F_p(x)\),
\[
 F_1'=1+(\eta-m_2)/4,\qquad F_2'=2F_1+(m_3-2c m_2+c\eta)/4.
\]
These identities continue \(F_2(X\sqrt z)\) holomorphically by integration of the moment expressions using \eqref{eq:original-7}.
In the slit sectors the \(z\) derivatives of \(m_2,m_3\) are bounded (arcsine integral and \(d\sigma/dz\) above), but \(d\eta/dz\sim-\pi^2 c_*^2/(2\delta\log^2(1/\delta))\).
Consequently
\begin{equation}
\frac{d^2}{dz^2}F_2(X\sqrt z)\sim
 \frac{4\pi^2}{(1-z)\log^2(1/(1-z))}.
 \label{eq:original-8}
\end{equation}

There are no other dominant circle singularities: in \eqref{eq:original-4} iteratively prune pendant tree edges in the original primal map, \emph{except} the root edge, which is protected.
The unique remaining rooted core (with no removable leaves off that edge) has \(m\ge1\) edges; each corner carries an arbitrary attached plane tree (including the empty choice).
Uniqueness also follows since the pruning keeps the cycles, the root and the paths needed to connect them, removing the attached unprotected trees.
Appending an open leaf keeps \(\ell\) unchanged, a closed one multiplies the weight by 2, by the displayed FK formula.
Rooted cores have no residual rotation symmetry.
Thus \(\sum_{n\ge1} Z_n(4)w^n=C(wT(w)^2)\), where \(T=1+3wT^2\) counts the attachments and \(C\) is a nonnegative power series counting the cores.
Since \(X^2=1/32<1/12\), \(T\) is regular, \(C\) is finite at \(X^2T(X^2)^2\) by \eqref{eq:original-4}, and all other points of \(|w|=X^2\) have strictly smaller \(|w T(w)^2|\), hence regular extensions.
This gives the indented-disk analyticity required in the Flajolet--Odlyzko singularity transfer theorem \citep{FlajoletOdlyzko}, together with \eqref{eq:original-7}.
Applying the logarithmic transfer (power exponent \(-1\)) to \eqref{eq:original-8} and then \eqref{eq:original-4} proves the first equivalent of \eqref{eq:original-3}.

\end{proof}

\subsection{Disk volume and the parity local limit theorem}\label{disk:volume}

\begin{lemma}[Disk-volume limit]\label{disk:volume-weak}

Write \(N_p\) for \(N\) in the degree \(p\) Boltzmann disk at \(X\) and \(v_p=\log^2 p/p^2\) for large \(p\).
Equations~\eqref{eq:original-6}--\eqref{eq:original-7} give
\begin{equation}
v_p N_p\ \Longrightarrow\ \frac{\pi^2}{8E},\qquad E\sim{\rm Exp}(1).
 \label{eq:original-9}
\end{equation}

\end{lemma}\begin{proof}

Indeed for \(x=X e^{-\lambda v_p}\), fixed \(\lambda>0\), \(p\alpha/c_*\to z_0=\pi\sqrt{\lambda/2}\).
In the moment integral in \(u'=c-y\in[\alpha,\beta]\), use \(v=p u'/c_*\); \((y/c_*)^p\to e^{-v}\) for bounded \(v\).
For \(v>z_0\),
\[
 \frac{p}{c_* \log p}\,r((v c_*/p)^2)\ \longrightarrow\
 \frac{\sqrt{v^2-z_0^2}}{\pi^2}.
\]
In fact the arcsine expectation of \(1/(U+v c_*/p)\) in \eqref{eq:original-6} is asymptotic to \(2\log p/(\pi\beta)\), by its sine-integral representation; insertion of \(v c_*/p\) in the denominator changes \(m_{-1}\) only by \(O(1)\), since \(\alpha\) has order \(1/p\).
The rescaled densities here are bounded by \(C v\) from \eqref{eq:original-6} and \(m_{-1}\), so exponential domination where \(y\ge0\) suffices; the region \(y<0\) contributes exponentially little.
Dividing by \(F_p(X)\), the limiting Laplace transform is \(\int_{z_0}^\infty e^{-v}\sqrt{v^2-z_0^2}\,dv\), equal to \(\int_0^\infty\exp(-u-z_0^2/(4u))du\) (change variable \(v=u+z_0^2/(4u)\), both branches, and integrate by parts).
This gives \eqref{eq:original-9}.

\end{proof}

\begin{lemma}[Uniform volume local limit theorem]\label{disk:volume-llt}

The lattice version needed for the finite-size split is
\begin{equation}
\mathbb P(N_p=j)=2v_p h(v_p j)+o(v_p),\qquad j\equiv p\pmod 2,\qquad
 h=d_{\pi/\sqrt8},
 \label{eq:original-10}
\end{equation}

uniformly over all integers $j$ of the indicated parity as $p\to\infty$ (extend \(h\) by zero to negative arguments). 

\end{lemma}\begin{proof}

Parity follows by counting triangle incidences.
Here are details to control Fourier inversion.
With \(\widetilde g_i=c_*^i g_i(X)\) one has

\[
 |f_{i+1}-f_i|\le C f_i/i,\qquad
 \widetilde g_i=2f_i+O(f_i/\sqrt i).
\]
The first bound follows using the moment formula at \(X\), the negative endpoint gap, \(t^i(1-t)\le (C/i)t^{i/2}\) on \([0,1]\), and the two-sided moment equivalent used above.
The second follows by representing \(\widetilde g_i/2\) as the expectation of \(f_j\) at \(j\) the sum of \(i\) iid geometric variables (mean one, distribution \(2^{-l-1}\) at \(l\ge0\)); use exponential deviations, first absolute fluctuation \(O(\sqrt i)\) and the first bound.

In \eqref{eq:original-5} follow always the larger-perimeter disk at a negative (bridge) step, with deterministic tie breaking, and the external remainder at a face step.
At current perimeter \(P=m\) large the next perimeter \(P'\) is \(m+d\); \(d=-i-2\) for \(0\le i\le(m-2)/2\), with probability \(c_*^{-2}w_d^- f_{m+d}/f_m\) where \(w_d^-=2 f_i\) (just \(f_i\) for the tie); or \(d=k-2\ge-1\) with probability \(c_*^{-2}w_d^+ f_{m+d}/f_m,\ w_d^+=\widetilde g_k\).
For fixed \(a>0\),
\[
 \mathbb E[(P')^{-a}\mid P=m]\le m^{-a}(1+C_a\log m/m).
\]
Indeed in the difference from \(m^{-a}\) one can drop \(d>m\) (negative terms).
In the remaining range \(1+d/m\) is bounded away from zero, \(f_{m+d}/f_m=1+O(|d|/m)\), and Taylor expansion gives linear contribution bounded relatively by \(C_a\log m/m\) and remainder by \(C_a m^{-2} O(m\log m)\).
Here \(\sum d(w_d^-+w_d^+)=O(\log m)\) on the indicated ranges (sum separately over each sign's stated indices), by \(\sum i|\widetilde g_i-2 f_i|<\infty\) and the log-over-square tail, which likewise gives \(\sum d^2(w_d^-+w_d^+)=O(m\log m)\).

Let \(\phi_p(\theta)=\mathbb E e^{i\theta N_p}\).
For sufficiently large \(K\), sufficiently small fixed \(\delta>0\), and \(K v_p\le\theta\le\delta\), take \(q>e\) solving \(\theta=\log^2 q/q^2\).
Put \(s=\sqrt{pq}\) and run at most \(\lfloor s/\log s\rfloor\) steps, stopping on crossing below \(s\); \(p\ge s\ge10q\) by choice of \(K\).
The probability of crossing is \(O_a((s/p)^a)\), by the stopped inverse-moment bound iterated (overshoot stays positive).
While above threshold, with probability at least \(c_1\log q/q\) a bridge step splits off an independent disk of perimeter \(i\in[q,2q]\).
By \eqref{eq:original-9}, \(\sup_{q\le i\le2q}|\phi_i(\theta)|\le1-c_2\) for all large \(q\), since the limiting characteristic function has modulus strictly below one on the relevant compact positive frequency interval.
Conditional on the peeling step types/perimeters, the swallowed disks are independent, contributing additively to the volume.
Hence by successive conditioning (drop the threshold survival indicator at each last step for the nonnegative product bound)
\[
 |\phi_p(\theta)|\le C_a(q/p)^{a/2}
        +(1-c_1 c_2\log q/q)^{\lfloor s/\log s\rfloor}
 \le C'_a(v_p/\theta)^{a/4}.
\]
We used that \(\sqrt{p/q}\log q/\log s\) grows faster than \(\log(p/q)\), and \(q/p\le\sqrt{v_p/\theta}\).
For \(\delta\le\theta\le\pi/2\) the same argument with fixed \(q\) and just perimeter \(i=2\) swallowed gives the last bound, since \(N_2\) can be 0 or 2.
This gives integrable domination on the rescaled frequency axis (take \(a>4\)); Fourier inversion on the parity lattice and \eqref{eq:original-9} now give \eqref{eq:original-10} by \(L^1\) convergence on that axis.

\end{proof}

\subsection{The two-disk cut and its counting bias}\label{disk:palm-proof}

\begin{proof}[Proof of the intensity assertion in Proposition~\ref{disk:palm}]

To mark a loop, cut out its triangle ring, leaving two disks with boundary degrees \(k,l\), colors blue and red.
This is the same ring decomposition as for gasket faces but now with two fillings.
Reroot at one of the \(k+l\) successive links distinguished in the ring (orientation specified by colors), replacing the directed FK root (with \(2n\) choices); counts have the corresponding ratio \((k+l)/(2n)\) since both markings kill sphere-map automorphisms.
Traversing from the link with blue side in a given orientation specifies the list interleaving both bases and hence disk boundary roots (if positive lengths).
Thus for \(k+l>0\),
\[
 \mathbb E\#\{\text{loops of type }(k,l)\}=
 \frac{4n}{k+l}\frac{\binom{k+l}{k}2^{-k-l} f_k f_l}{X^{2n}Z_n(4)}
 \ \mathbb P(N_k+N_l=2n-k-l),
\]
where the fillings are independent before total conditioning.
Also \(u-N_k/(2n)\) is between zero and \((k+l)/(2n)\), since each Tutte triangle has exactly two flags.
For \(k/B_n\) in compact subsets of \((0,\infty)\), summation over \(l\) is exponentially concentrated at \(l/k=1+o(1)\) by the binomial kernel of total mass 2 (tails outside e.g. \(k\pm k^{3/4}\) dominate no polynomial factor; \(f_j\) are bounded).
For $k\sim sB_n$, the perimeter and area units satisfy
\[
 v_k\sim\frac{\pi^2}{16s^2n},\qquad
 \frac{N_k}{2n}\Longrightarrow\frac{s^2}{E},\qquad
 f_k,f_l\sim\frac1{s^2n\log n}.
\]
The span-two lattice in Lemma~\ref{disk:volume-llt} therefore gives
\[
 \Prob(N_k=j)=\frac1n\bigl(d_s(j/(2n))+o(1)\bigr),
 \qquad j\equiv k\pmod2,
\]
uniformly in $j$ and for $s$ in the stated compact range.  Thus the area
mesh for one filling is $1/n$.  This lattice factor is distinct from the
loop fugacity $2$, the rerooting factor $2n/(k+l)$, and the total mass $2$
of the ring kernel summed over $l$.
The total constraint leaves the first fraction on a grid of spacing \(1/n\) with the complementary parity automatically satisfied.
Its constrained probability measure multiplied by \(n\) thus tends weakly to \(d_s(v)d_s(1-v)\,dv\) along every convergent \(s\), with the same limit using \(u\).
Inserting and then summing with \(s\)-spacing \(1/B_n\) proves the second formula of \eqref{eq:original-3}.

\end{proof}

These volume and first-intensity limits do not by themselves give a coupling of the entire maps to the critical surface.
We next build the cutting and peeling comparisons needed to continue.
